\section{Banco de questões — Raciocínio Lógico}

\subsection*{N1 — Fundamento competitivo}
\begin{questao}{\nivel{N1} 03.01 — Cebraspe/Certo ou Errado}A negação de $p\to q$ é logicamente equivalente a $p\land\neg q$.\end{questao}
\begin{questao}{\nivel{N1} 03.02 — Cebraspe/Certo ou Errado}A proposição $p\to q$ é falsa exatamente quando $p$ é verdadeira e $q$ é falsa.\end{questao}
\begin{questao}{\nivel{N1} 03.03 — Cebraspe/Certo ou Errado}$p\to q$ é equivalente a $q\to p$.\end{questao}
\begin{questao}{\nivel{N1} 03.04 — Cebraspe/Certo ou Errado}$p\to q$ é equivalente a $\neg q\to\neg p$.\end{questao}
\begin{questao}{\nivel{N1} 03.05 — Cebraspe/Certo ou Errado}A negação de ``Todos os contratos foram revisados'' é ``Nenhum contrato foi revisado''.\end{questao}
\begin{questao}{\nivel{N1} 03.06 — Cebraspe/Certo ou Errado}$p\lor\neg p$ é tautologia.\end{questao}
\begin{questao}{\nivel{N1} 03.07 — Cebraspe/Certo ou Errado}$p\land\neg p$ é contradição.\end{questao}
\begin{questao}{\nivel{N1} 03.08 — Cebraspe/Certo ou Errado}Se $A\subseteq B$, então todo elemento de $A$ pertence a $B$.\end{questao}
\begin{questao}{\nivel{N1} 03.09 — Cebraspe/Certo ou Errado}Aumentar um valor em 20\% e depois reduzi-lo em 20\% retorna ao valor inicial.\end{questao}
\begin{questao}{\nivel{N1} 03.10 — Cebraspe/Certo ou Errado}Se duas grandezas são inversamente proporcionais, o produto entre valores correspondentes permanece constante, sob as hipóteses do modelo.\end{questao}

\subsection*{N2 — Aplicação}
\begin{questao}{\nivel{N2} 03.11 — FGV/condicional}A proposição ``Se o relatório foi aprovado, então foi assinado'' é falsa apenas quando:\begin{enumerate}[label=\Alph*)]\item o relatório não foi aprovado e não foi assinado;\item foi aprovado e não foi assinado;\item não foi aprovado e foi assinado;\item foi aprovado e assinado;\item não foi assinado, independentemente da aprovação.\end{enumerate}\end{questao}
\begin{questao}{\nivel{N2} 03.12 — FCC/equivalência}Uma forma equivalente a ``Se há falha crítica, então o serviço é interrompido'' é:\begin{enumerate}[label=\Alph*)]\item Se o serviço é interrompido, há falha crítica.\item Se o serviço não é interrompido, então não há falha crítica.\item Há falha crítica e o serviço não é interrompido.\item Se não há falha crítica, o serviço não é interrompido.\item Há falha crítica se, e somente se, o serviço é interrompido.\end{enumerate}\end{questao}
\begin{questao}{\nivel{N2} 03.13 — Cebraspe/Certo ou Errado}A negação de ``O sistema é seguro e disponível'' é ``O sistema não é seguro ou não é disponível''.\end{questao}
\begin{questao}{\nivel{N2} 03.14 — FGV/De Morgan}A negação de ``O contrato tem SLA ou possui cláusula de penalidade'' é:\begin{enumerate}[label=\Alph*)]\item não tem SLA ou não possui penalidade;\item não tem SLA e não possui penalidade;\item tem SLA e possui penalidade;\item se tem SLA, não possui penalidade;\item tem exatamente uma das duas condições.\end{enumerate}\end{questao}
\begin{questao}{\nivel{N2} 03.15 — FCC/argumento}Premissas: ``Se há acesso indevido, então há risco.'' ``Há acesso indevido.'' Conclusão válida:\begin{enumerate}[label=\Alph*)]\item não há risco;\item há risco;\item não há acesso indevido;\item risco implica acesso indevido;\item nada pode ser concluído.\end{enumerate}\end{questao}
\begin{questao}{\nivel{N2} 03.16 — Cebraspe/Certo ou Errado}Do argumento $p\to q$ e $q$, concluir $p$ constitui, em geral, a falácia da afirmação do consequente.\end{questao}
\begin{questao}{\nivel{N2} 03.17 — FGV/quantificadores}A negação de ``Todo auditor analisou algum processo'' é:\begin{enumerate}[label=\Alph*)]\item Nenhum auditor analisou processo algum.\item Existe auditor que não analisou processo algum.\item Todo auditor deixou de analisar algum processo.\item Existe processo não analisado por nenhum auditor.\item Algum auditor analisou todos os processos.\end{enumerate}\end{questao}
\begin{questao}{\nivel{N2} 03.18 — FCC/conjuntos}Em um grupo, 70 pessoas usam sistema A, 50 usam B e 30 usam ambos. Quantas usam A ou B?\begin{enumerate}[label=\Alph*)]\item 60;\item 90;\item 100;\item 120;\item 150.\end{enumerate}\end{questao}
\begin{questao}{\nivel{N2} 03.19 — FGV/conjuntos}Em 100 servidores, 60 dominam Excel, 45 dominam Python e 25 dominam ambos. Quantos não dominam nenhum dos dois?\begin{enumerate}[label=\Alph*)]\item 10;\item 20;\item 25;\item 30;\item 40.\end{enumerate}\end{questao}
\begin{questao}{\nivel{N2} 03.20 — Cebraspe/Certo ou Errado}Se 40\% de 250 processos apresentam determinada característica, então 100 processos a apresentam.\end{questao}
\begin{questao}{\nivel{N2} 03.21 — FCC/porcentagem}Um valor de R\$ 800 mil aumenta 15\%. O novo valor é:\begin{enumerate}[label=\Alph*)]\item R\$ 880 mil;\item R\$ 900 mil;\item R\$ 920 mil;\item R\$ 940 mil;\item R\$ 960 mil.\end{enumerate}\end{questao}
\begin{questao}{\nivel{N2} 03.22 — FGV/percentuais sucessivos}Um preço aumenta 10\% e depois 20\%. O aumento total é:\begin{enumerate}[label=\Alph*)]\item 28\%;\item 30\%;\item 32\%;\item 33\%;\item 35\%.\end{enumerate}\end{questao}
\begin{questao}{\nivel{N2} 03.23 — Cebraspe/Certo ou Errado}Uma redução de 25\% seguida de aumento de 25\% produz valor final inferior ao inicial.\end{questao}
\begin{questao}{\nivel{N2} 03.24 — FCC/proporção}Se 6 analistas revisam 180 processos em 5 dias, mantendo produtividade constante, 10 analistas revisariam 300 processos em:\begin{enumerate}[label=\Alph*)]\item 3 dias;\item 4 dias;\item 5 dias;\item 6 dias;\item 10 dias.\end{enumerate}\end{questao}
\begin{questao}{\nivel{N2} 03.25 — FGV/equação}A soma de um número com seu dobro é 45. O número é:\begin{enumerate}[label=\Alph*)]\item 10;\item 12;\item 15;\item 20;\item 30.\end{enumerate}\end{questao}

\subsection*{N3 — Integração de concurso}
\begin{questao}{\nivel{N3} 03.26 — FGV/assertivas}Considere: I. $p\to q\equiv\neg p\lor q$. II. $\neg(p\land q)\equiv\neg p\land\neg q$. III. $p\leftrightarrow q\equiv(p\to q)\land(q\to p)$. Assinale:\begin{enumerate}[label=\Alph*)]\item apenas I;\item apenas II;\item I e III;\item II e III;\item I, II e III.\end{enumerate}\end{questao}
\begin{questao}{\nivel{N3} 03.27 — Cebraspe/Certo ou Errado}A negação de ``Se o sistema falha, o alerta é emitido'' pode ser expressa por ``O sistema falha e o alerta não é emitido''.\end{questao}
\begin{questao}{\nivel{N3} 03.28 — FCC/necessária-suficiente}Na proposição ``Se o contrato está vigente, então o pagamento pode ocorrer'', a vigência é condição:\begin{enumerate}[label=\Alph*)]\item necessária para a implicação, e pagamento é suficiente;\item suficiente para a conclusão expressa e o consequente é condição necessária para o antecedente;\item necessária e suficiente;\item sem relação lógica;\item equivalente ao consequente.\end{enumerate}\end{questao}
\begin{questao}{\nivel{N3} 03.29 — FGV/argumento}Premissas: $p\to q$, $q\to r$, $\neg r$. Conclusão necessária:\begin{enumerate}[label=\Alph*)]\item $p$;\item $q$;\item $\neg p$;\item $r$;\item $p\land q$.\end{enumerate}\end{questao}
\begin{questao}{\nivel{N3} 03.30 — Cebraspe/Certo ou Errado}Das premissas ``Todo A é B'' e ``Algum B é C'', não se conclui necessariamente que ``Algum A é C''.\end{questao}
\begin{questao}{\nivel{N3} 03.31 — FCC/diagramas}Se todo auditor é servidor e nenhum terceirizado é servidor, então é necessariamente verdadeiro que:\begin{enumerate}[label=\Alph*)]\item nenhum auditor é terceirizado;\item todo terceirizado é auditor;\item algum auditor é terceirizado;\item todo servidor é auditor;\item nenhum servidor é terceirizado e nenhum terceirizado existe.\end{enumerate}\end{questao}
\begin{questao}{\nivel{N3} 03.32 — FGV/conjuntos três}Em 200 pessoas: $|A|=120$, $|B|=90$, $|C|=80$, $|A\cap B|=50$, $|A\cap C|=40$, $|B\cap C|=30$ e $|A\cap B\cap C|=20$. Quantas pertencem a pelo menos um conjunto?\begin{enumerate}[label=\Alph*)]\item 150;\item 170;\item 180;\item 190;\item 200.\end{enumerate}\end{questao}
\begin{questao}{\nivel{N3} 03.33 — Cebraspe/Certo ou Errado}No cenário da questão anterior, 30 pessoas pertencem a $A\cap B$ mas não a $C$.\end{questao}
\begin{questao}{\nivel{N3} 03.34 — FCC/porcentagem inversa}Após desconto de 20\%, um produto custa R\$ 800. O preço original era:\begin{enumerate}[label=\Alph*)]\item R\$ 900;\item R\$ 960;\item R\$ 1.000;\item R\$ 1.020;\item R\$ 1.200.\end{enumerate}\end{questao}
\begin{questao}{\nivel{N3} 03.35 — FGV/variação}Uma taxa cai de 25\% para 20\%. A redução foi de:\begin{enumerate}[label=\Alph*)]\item 5\% relativa;\item 5 pontos percentuais e 20\% relativa;\item 20 pontos percentuais;\item 25\% relativa;\item 80\% relativa.\end{enumerate}\end{questao}
\begin{questao}{\nivel{N3} 03.36 — Cebraspe/Certo ou Errado}A diferença entre 30\% e 20\% é 10 pontos percentuais, o que corresponde a redução relativa de um terço quando se parte de 30\%.\end{questao}
\begin{questao}{\nivel{N3} 03.37 — FCC/proporção inversa}Oito servidores realizam tarefa em 15 dias. Mantida produtividade, 12 servidores realizam a mesma tarefa em:\begin{enumerate}[label=\Alph*)]\item 8 dias;\item 10 dias;\item 12 dias;\item 15 dias;\item 22,5 dias.\end{enumerate}\end{questao}
\begin{questao}{\nivel{N3} 03.38 — FGV/sistema}Em uma equipe há auditores e analistas, total 30. O número de analistas é o dobro do de auditores. Quantos auditores há?\begin{enumerate}[label=\Alph*)]\item 5;\item 10;\item 15;\item 20;\item 25.\end{enumerate}\end{questao}
\begin{questao}{\nivel{N3} 03.39 — Cebraspe/Certo ou Errado}Se a média de 5 valores é 20, a soma dos valores é 100.\end{questao}
\begin{questao}{\nivel{N3} 03.40 — FCC/problema}Um órgão analisou 60\% de 500 processos. Dos analisados, 10\% tinham falha crítica. Quantos processos analisados tinham falha crítica?\begin{enumerate}[label=\Alph*)]\item 20;\item 25;\item 30;\item 50;\item 60.\end{enumerate}\end{questao}

\subsection*{N4 — Avançado}
\begin{questao}{\nivel{N4} 03.41 — Cebraspe/caso integrado}Considere: ``Se o backup é válido, então a restauração é possível.'' Foi constatado que a restauração não é possível. É logicamente válido concluir que o backup não é válido.\end{questao}
\begin{questao}{\nivel{N4} 03.42 — FGV/falácia}Premissa: ``Se há fraude, há indício documental.'' Constatou-se indício documental. Concluir que houve fraude comete:\begin{enumerate}[label=\Alph*)]\item modus ponens;\item modus tollens;\item afirmação do consequente;\item negação do consequente;\item contraposição válida.\end{enumerate}\end{questao}
\begin{questao}{\nivel{N4} 03.43 — FCC/quantificador}A negação de ``Existe servidor que revisou todos os processos'' é:\begin{enumerate}[label=\Alph*)]\item Nenhum servidor revisou processo algum.\item Para todo servidor, existe ao menos um processo que ele não revisou.\item Existe processo não revisado por nenhum servidor.\item Todo servidor revisou algum processo.\item Nenhum processo foi revisado por todos.\end{enumerate}\end{questao}
\begin{questao}{\nivel{N4} 03.44 — Cebraspe/Certo ou Errado}A negação de $\exists x\forall y\,P(x,y)$ é $\forall x\exists y\,\neg P(x,y)$.\end{questao}
\begin{questao}{\nivel{N4} 03.45 — FGV/conjuntos}Em uma população de 300 pessoas, 180 usam A, 160 usam B e 120 usam ambos. O número que usa exatamente um dos dois é:\begin{enumerate}[label=\Alph*)]\item 80;\item 100;\item 120;\item 140;\item 220.\end{enumerate}\end{questao}
\begin{questao}{\nivel{N4} 03.46 — FCC/percentual composto}Um orçamento sofre redução de 10\%, depois aumento de 10\% e, por fim, aumento de 5\%. Em relação ao inicial, o resultado é aproximadamente:\begin{enumerate}[label=\Alph*)]\item -5,0\%;\item -1,0\%;\item +3,95\%;\item +5,0\%;\item +15,0\%.\end{enumerate}\end{questao}
\begin{questao}{\nivel{N4} 03.47 — Cebraspe/Certo ou Errado}Se uma quantidade aumenta de 80 para 100, o aumento relativo é 25\%; se depois volta de 100 para 80, a redução relativa é 20\%.\end{questao}
\begin{questao}{\nivel{N4} 03.48 — FGV/alocação}Três equipes A, B e C recebem processos na razão 2:3:5. Se há 400 processos, a equipe C recebe:\begin{enumerate}[label=\Alph*)]\item 80;\item 120;\item 160;\item 200;\item 250.\end{enumerate}\end{questao}
\begin{questao}{\nivel{N4} 03.49 — FCC/equação aplicada}Uma contratação custa taxa fixa de R\$ 50 mil mais R\$ 2 mil por unidade. Se o total foi R\$ 170 mil, o número de unidades foi:\begin{enumerate}[label=\Alph*)]\item 50;\item 60;\item 70;\item 85;\item 110.\end{enumerate}\end{questao}
\begin{questao}{\nivel{N4} 03.50 — FGV/lógica e linguagem}A frase ``O pagamento ocorrerá somente se houver aceite'' estabelece que:\begin{enumerate}[label=\Alph*)]\item aceite é condição suficiente para pagamento;\item aceite é condição necessária para pagamento;\item pagamento é condição necessária para aceite;\item pagamento e aceite são equivalentes;\item aceite é irrelevante.\end{enumerate}\end{questao}
